\documentclass{article}
\usepackage{graphicx} % Required for inserting images
\usepackage[onefile]{leantex}
\usepackage[margin=0.5in]{geometry}
\title{Lean Test No. 4: 161 Assignment 4 part 2}
\author{Ryland Gross}
\date{February 2026}

\begin{document}

\maketitle

\section{Imports, Variables and Options}

\begin{lean}
import Mathlib.Tactic
import Mathlib.Util.Delaborators
import Mathlib.Data.Real.Basic

set_option warningAsError false


variable {α β : Type} (p q : α → Prop) (r : α → β → Prop)
\end{lean}

\section{Exercise 1}

\begin{lean}
example : (∀ x, p x) ∧ (∀ x, q x) → ∀ x, p x ∧ q x := by
  intros h x
  constructor
  · exact h.left x
  exact h.right x

example : (∀ x, p x) ∨ (∀ x, q x) → ∀ x, p x ∨ q x := by
  intros h x
  rcases h with h1 | h2
  · left
    exact h1 x
  right
  exact h2 x

example : (∃ x, ∀ y, r x y) → ∀ y, ∃ x, r x y := by
  intros h y
  rcases h with ⟨x, h11⟩
  use x
  exact h11 y
\end{lean}

\section{Some Definitions and Lemmas}

\begin{lean}
def approaches_at (f : ℝ → ℝ) (b : ℝ) (a : ℝ) :=
∀ ε > 0, ∃ δ > 0, ∀ x, abs (x - a) < δ → abs (f x - b) < ε

def continuous (f : ℝ → ℝ) := ∀ x, approaches_at f (f x) x

theorem my_continuous_add_right2 {f : ℝ → ℝ} (ctsf : continuous f) (r : ℝ) :
  continuous (fun x ↦ f x + r) := by
  intros x ε h
  have h1 := ctsf x ε h
  rcases h1 with ⟨δ, h11⟩
  use δ
  constructor
  · exact h11.left
  intros x_1 h
  rw[add_sub_add_right_eq_sub]
  exact h11.right x_1 h


-- Since `f x - r` is the same as `f x + (- r)`, the following is an instance
-- of the previous theorem.
theorem continuous_sub {f : ℝ → ℝ} (ctsf : continuous f) (r : ℝ) :
  continuous (fun x ↦ f x - r) :=
  my_continuous_add_right2 ctsf (-r)
\end{lean}

\section{Exercise 4 Test}

Here is the intermediate value theorem, we will solve the following problem modulo it
\begin{lean}
theorem ivt {f : ℝ → ℝ} {a b : ℝ} (aleb : a ≤ b)
    (ctsf : continuous f) (hfa : f a < 0) (hfb : 0 < f b) :
  ∃ x, a ≤ x ∧ x ≤ b ∧ f x = 0 :=
sorry
\end{lean}

Now, we actually exhibit a more genral version,
\begin{lean}
theorem ivtnolinarith {f : ℝ → ℝ} {a b c : ℝ} (aleb : a ≤ b)
    (ctsf : continuous f) (hfa : f a < c) (hfb : c < f b) :
  ∃ x, a ≤ x ∧ x ≤ b ∧ f x = c := by
  unfold continuous at ctsf
  unfold approaches_at at ctsf
  let g := fun x ↦ f x - c
  have h1 : ∃ x, a ≤ x ∧ x ≤ b ∧ g x = 0 := by
    have ctsg : continuous g := by
      apply continuous_sub ctsf
    apply ivt aleb ctsg
    · exact sub_neg_of_lt hfa
    exact sub_pos_of_lt hfb
  dsimp [g] at h1
  rcases h1 with ⟨x, h11⟩
  use x
  constructor
  · exact h11.left
  constructor
  · exact h11.right.left
  exact eq_of_sub_eq_zero h11.right.right
\end{lean}

Now we are done. \textbf{Q.E.D.}

\end{document}
